\mxunit{seconde-s02}{Vecteurs du plan}\mxsubtitle{14--30 septembre 2026}
\section{Translation et vecteurs}
\begin{keypoint}[Définition]
La translation qui transforme $A$ en $B$ associe à tout point $M$ l'unique point $M'$ tel que $ABM'M$ soit un parallélogramme, éventuellement aplati. Elle est appelée translation de vecteur $\vect{AB}$.
Un vecteur non nul possède une direction, un sens et une norme. $\|\vect{AB}\|=AB$. Le vecteur $\vect{AA}$ est le vecteur nul, noté $\vec0$ ; il n'a ni direction ni sens.
\end{keypoint}
\begin{center}\begin{tikzpicture}[>=Stealth,scale=.9]
\coordinate(A) at (0,0);\coordinate(B) at (4,1);\coordinate(M) at (1,2);\coordinate(N) at (5,3);
\draw[dashed](A)--(M)--(N)--(B);\draw[->,thick](A)--(B);\draw[->,thick](M)--(N);
\foreach \p/\n in {A/A,B/B,M/M,N/M'}\fill(\p)circle(1.5pt) node[above left]{$\n$};
\end{tikzpicture}\end{center}
\begin{keypoint}[Vecteurs égaux et opposés]
$\vect{AB}=\vect{CD}$ lorsque la translation de vecteur $\vect{AB}$ transforme $C$ en $D$. Deux vecteurs non nuls sont égaux si et seulement s'ils ont même direction, même sens et même norme.
Le vecteur opposé de $\vect{AB}$ est $\vect{BA}$ : $\vect{BA}=-\vect{AB}$.
\end{keypoint}
\begin{example}Tracer un segment $[AB]$ et placer un point $C$ hors de $(AB)$. Construire $D$ tel que $\vect{AB}=\vect{CD}$. Nommer le parallélogramme obtenu. Construire $E$ tel que $\vect{CE}=-\vect{AB}$.\end{example}
\mxanswerspace[4]{$D$ se construit en reportant depuis $C$ la direction, le sens et la longueur de $\vect{AB}$. Le parallélogramme est $ABDC$, dans cet ordre. $E$ est le symétrique de $D$ par rapport à $C$, car $\vect{CE}=-\vect{CD}$.}

\mxsection{Somme, différence et relation de Chasles}
\begin{keypoint}[Somme de vecteurs]
La somme $\vec u+\vec v$ correspond à l'enchaînement des translations de vecteurs $\vec u$ et $\vec v$. Pour tous points $A,B,C$ :
\[\boxed{\vect{AB}+\vect{BC}=\vect{AC}}\qquad\text{(relation de Chasles).}\]
On a $\vec u-\vec v=\vec u+(-\vec v)$ et $\vect{AB}+\vect{BA}=\vec0$.
\end{keypoint}
\begin{center}\begin{tikzpicture}[>=Stealth,scale=1.0]
\draw[->,thick](0,0)--(3,0)node[midway,below]{$\vec u$};
\draw[->,thick](3,0)--(4,2)node[midway,right]{$\vec v$};
\draw[->,very thick](0,0)--(4,2)node[midway,above left]{$\vec u+\vec v$};
\draw[dashed](0,0)--(1,2)--(4,2);
\node[below left]at(0,0){$A$};\node[below]at(3,0){$B$};\node[above right]at(4,2){$C$};
\end{tikzpicture}\end{center}
\begin{keypoint}[Règle du parallélogramme]
$ABDC$ est un parallélogramme, éventuellement aplati, si et seulement si $\vect{AD}=\vect{AB}+\vect{AC}$.
\end{keypoint}
\begin{example}Simplifier : $\vect{AB}+\vect{BC}$ ; $\vect{AB}+\vect{CA}$ ; $\vect{AB}+\vect{BC}+\vect{CD}$ ; $\vect{AB}-\vect{AC}$.\end{example}
\mxanswerspace[3.2]{$\vect{AC}$ ; $\vect{CA}+\vect{AB}=\vect{CB}$ ; $\vect{AD}$ ; $\vect{AB}+\vect{CA}=\vect{CB}$. La commutativité de l'addition permet de réordonner les vecteurs.}
\begin{example}Dans un parallélogramme $ABCD$, justifier que $\vect{AB}+\vect{AD}=\vect{AC}$.\end{example}
\mxanswerspace[2.4]{Dans $ABCD$, $\vect{AD}=\vect{BC}$. Donc $\vect{AB}+\vect{AD}=\vect{AB}+\vect{BC}=\vect{AC}$.}

\mxsection{Produit d'un vecteur par un réel}
\begin{keypoint}[Définition et règles de calcul]
Pour $k\in\R$, $k\vec u=\vec0$ si $k=0$ ou $\vec u=\vec0$. Sinon, $k\vec u$ a la même direction que $\vec u$, le même sens si $k>0$, le sens opposé si $k<0$, et sa norme vaut $|k|\|\vec u\|$.
\[(k+\ell)\vec u=k\vec u+\ell\vec u,\quad k(\vec u+\vec v)=k\vec u+k\vec v,\quad k(\ell\vec u)=(k\ell)\vec u.\]
\end{keypoint}
\begin{example}Pour un vecteur $\vec u$ de norme $4$ cm, représenter $2\vec u$, $-\vec u$ et $\frac12\vec u$ à partir d'une même origine. Préciser leurs normes.\end{example}
\mxanswerspace[3.2]{Les normes sont $8$ cm, $4$ cm, $2$ cm. Les vecteurs $2\vec u$ et $\frac12\vec u$ ont le sens de $\vec u$ ; $-\vec u$ a le sens opposé.}
\begin{keypoint}[Caractérisations du milieu]
$I$ est le milieu de $[AB]$ si et seulement si $\vect{AI}=\vect{IB}$, ou encore $\vect{AI}=\frac12\vect{AB}$, ou $\vect{IA}+\vect{IB}=\vec0$.
\end{keypoint}
\begin{example}Dans un triangle $ABC$, $I$ est le milieu de $[AB]$ et $J$ celui de $[AC]$. Montrer que $\vect{IJ}=\frac12\vect{BC}$.\end{example}
\mxanswerspace[4.4]{Par Chasles, $\vect{IJ}=\vect{IA}+\vect{AJ}=-\frac12\vect{AB}+\frac12\vect{AC}=\frac12(\vect{BA}+\vect{AC})=\frac12\vect{BC}$. On retrouve le parallélisme des droites $(IJ)$ et $(BC)$ lorsque le triangle n'est pas aplati.}
\begin{fluency}[5--10 min]Réduire $3\vec u-5\vec u$ ; calculer $(-3)^2$ ; résoudre $2x+1=7$.\mxanswer{$-2\vec u$ ; $9$ ; $x=3$.}\end{fluency}

\mxsection{Repères, bases et coordonnées}
\begin{keypoint}[Définitions]
Deux vecteurs non colinéaires $\vec i,\vec j$ forment une base du plan. Tout vecteur $\vec u$ s'écrit de manière unique $\vec u=x\vec i+y\vec j$ ; ses coordonnées sont $\pt{x;y}$ ou $\vc{x}{y}$.
Un repère $\pt{O;\vec i,\vec j}$ associe au point $M$ les coordonnées de $\vect{OM}$. Il est \textbf{orthonormé} si les vecteurs de base sont perpendiculaires et de norme $1$.
\end{keypoint}
\begin{center}\begin{tikzpicture}[>=Stealth,scale=.8]
\draw[step=1,black!20](-1,-1)grid(6,4);\draw[->](-1,0)--(6.3,0)node[right]{$x$};\draw[->](0,-1)--(0,4.3)node[above]{$y$};
\draw[->,very thick](0,0)--(1,0)node[midway,below]{$\vec i$};\draw[->,very thick](0,0)--(0,1)node[midway,left]{$\vec j$};
\draw[->,thick](1,1)--(4,3)node[midway,above left]{$\vect{AB}$};\draw[dashed](1,1)--(4,1)--(4,3);
\node[below]at(1,1){$A$};\node[above]at(4,3){$B$};\node[below left]at(0,0){$O$};
\end{tikzpicture}\end{center}
Ici $\vect{AB}=3\vec i+2\vec j$ : on se déplace de trois unités vers la droite et deux vers le haut.
\begin{keypoint}[Calculs avec les coordonnées]
Si $A\pt{x_A;y_A}$, $B\pt{x_B;y_B}$, $\vec u\pt{x;y}$ et $\vec v\pt{x';y'}$ :
\[\vect{AB}=\vc{x_B-x_A}{y_B-y_A},\quad \vec u+\vec v=\vc{x+x'}{y+y'},\quad k\vec u=\vc{kx}{ky}.\]
Deux vecteurs sont égaux si et seulement si leurs coordonnées respectives sont égales.
\end{keypoint}
\begin{example}Pour $\vec u\pt{3;2}$ et $\vec v\pt{-1;4}$, calculer $\vec u+\vec v$ et $2\vec u-\vec v$. Écrire $\vec w\pt{5;-3}$ dans la base $\pt{\vec i;\vec j}$.\end{example}
\mxanswerspace[2.4]{$\vec u+\vec v=\pt{2;6}$ ; $2\vec u-\vec v=\pt{7;0}$ ; $\vec w=5\vec i-3\vec j$.}

\mxsection{Norme, distance et milieu}
\begin{keypoint}[Propriétés]
Dans un repère \textbf{orthonormé}, pour $\vec u\pt{x;y}$ et deux points $A,B$ :
\[\|\vec u\|=\sqrt{x^2+y^2},\qquad AB=\sqrt{(x_B-x_A)^2+(y_B-y_A)^2}.\]
Dans tout repère, le milieu de $[AB]$ est $I\left(\frac{x_A+x_B}{2};\frac{y_A+y_B}{2}\right)$.
\end{keypoint}
\mxwriteproof[3]{La formule de distance découle du théorème de Pythagore appliqué aux déplacements horizontal et vertical. Pour le milieu, $\vect{AI}=\frac12\vect{AB}$ donne $x_I-x_A=\frac{x_B-x_A}{2}$ et $y_I-y_A=\frac{y_B-y_A}{2}$ ; on isole $x_I,y_I$.}
\begin{example}Dans un repère orthonormé, $A\pt{3;0}$, $B\pt{5;3}$, $C\pt{4;0}$ et $D\pt{6;4}$.
\begin{enumerate}
\item Calculer $\vect{AB}$, $\vect{CD}$, puis $3\vect{AB}-5\vect{CD}$.
\item Déterminer $U$ tel que $\vect{BU}=2\vect{AB}-\vect{CD}$.
\item Déterminer $E$ tel que $ABCE$ soit un parallélogramme, puis son centre $I$.
\item Calculer $AB$ et $CD$.
\end{enumerate}\end{example}
\mxanswerspace[6]{1. $\vect{AB}=\pt{2;3}$ et $\vect{CD}=\pt{2;4}$, donc $3\vect{AB}-5\vect{CD}=\pt{-4;-11}$.\\
2. $\vect{BU}=\pt{2;2}$, donc $U=B+\pt{2;2}=\pt{7;5}$.\\
3. $ABCE$ est un parallélogramme si et seulement si $\vect{EC}=\vect{AB}$, soit $\pt{4-x_E;-y_E}=\pt{2;3}$, donc $E\pt{2;-3}$. Le centre est le milieu de $[AC]$ : $I\pt{7/2;0}$. Le milieu de $[BE]$ donne le même résultat.\\
4. $AB=\sqrt{13}$ et $CD=\sqrt{20}=2\sqrt5$.}

\mxsection{Colinéarité et déterminant}
\begin{keypoint}[Définition]
Deux vecteurs non nuls sont colinéaires s'ils ont la même direction. Le vecteur nul est colinéaire à tout vecteur. Si $\vec v\ne\vec0$, $\vec u$ et $\vec v$ sont colinéaires si et seulement s'il existe un réel $k$ tel que $\vec u=k\vec v$.
\end{keypoint}
\begin{pitfall}Dans la dernière caractérisation, la condition $\vec v\ne\vec0$ est nécessaire : un vecteur non nul ne peut pas être un multiple du vecteur nul.\end{pitfall}
\begin{keypoint}[Déterminant et critère de colinéarité]
Pour $\vec u\pt{x;y}$, $\vec v\pt{x';y'}$, $\det(\vec u,\vec v)=xy'-yx'$.
\[\vec u,\vec v\text{ colinéaires}\quad\Longleftrightarrow\quad\det(\vec u,\vec v)=0.\]
\end{keypoint}
\mxwriteproof[4.2]{Si $\vec u=k\vec v$, alors $xy'-yx'=kx'y'-ky'x'=0$. Le cas où un vecteur est nul est immédiat.
Réciproquement, supposons $\vec v\ne\vec0$ et $xy'-yx'=0$. Si $x'\ne0$, posons $k=x/x'$. L'égalité donne $y=ky'$, donc $\vec u=k\vec v$. Si $x'=0$, alors $y'\ne0$, d'où $x=0$ ; avec $k=y/y'$, on obtient encore $\vec u=k\vec v$.}
\begin{example}Étudier la colinéarité : (a) $\vec u\pt{3;2}$ et $\vec v\pt{-6;-5}$ ; (b) $\vec w\pt{4;1}$ et $\vec t\pt{3;3/4}$. Si $\vect{AB}=2\vect{CD}$ et $\vect{BD}=-5\vect{CD}$, relier $\vect{AB}$ et $\vect{BD}$.\end{example}
\mxanswerspace[3.2]{(a) $\det(\vec u,\vec v)=-15+12=-3\ne0$ : non colinéaires. (b) $\vec t=\frac34\vec w$ : colinéaires ; le déterminant vaut $4\times\frac34-3=0$. Enfin $\vect{AB}=-\frac25\vect{BD}$, y compris si les vecteurs sont nuls.}

\mxsection{Alignement et parallélisme}
\begin{keypoint}[Applications]
Pour $A\ne B$ et $C\ne D$, les droites $(AB)$ et $(CD)$ sont parallèles (éventuellement confondues) si et seulement si $\det(\vect{AB},\vect{CD})=0$.
Les points $A,B,C$ sont alignés si et seulement si $\det(\vect{AB},\vect{AC})=0$.
\end{keypoint}
\begin{example}Montrer que $(AB)$ et $(CD)$ sont parallèles avec $A\pt{2;-4}$, $B\pt{10;-2}$, $C\pt{-16;-4}$ et $D\pt{12;3}$.\end{example}
\mxanswerspace[3.2]{$\vect{AB}=\pt{8;2}$ et $\vect{CD}=\pt{28;7}$. Leur déterminant vaut $8\times7-2\times28=0$. Les points de chaque paire sont distincts : les deux droites existent et sont parallèles.}
\begin{example}Montrer que $E\pt{0;1/2}$, $F\pt{2;1}$ et $G\pt{6;2}$ sont alignés.\end{example}
\mxanswerspace[3.2]{$\vect{EF}=\pt{2;1/2}$ et $\vect{EG}=\pt{6;3/2}=3\vect{EF}$. Le déterminant vaut $2\times3/2-(1/2)\times6=0$ ; donc les points sont alignés.}
\begin{example}Avec $A\pt{-1;1}$, $B\pt{3;2}$, $C\pt{-2;-3}$, $D\pt{6;-1}$, $E\pt{5;0}$, montrer que $(AB)\parallel(CD)$ et que $E,B,D$ sont alignés.\end{example}
\mxanswerspace[3.2]{$\vect{AB}=\pt{4;1}$ et $\vect{CD}=\pt{8;2}=2\vect{AB}$. Puis $\vect{EB}=\pt{-2;2}$ et $\vect{ED}=\pt{1;-1}=-\frac12\vect{EB}$. Les deux résultats suivent.}
